\chapter{Design}

% TODO: Find citations

As decisões de design do sistema foram fundamentadas tanto na análise dos
aplicativos similares quanto nos resultados do questionário aplicado durante a
etapa de pesquisa. Esses elementos permitiram identificar padrões recorrentes
de interface e compreender as preferências do público-alvo, orientando a
definição da identidade visual, da linguagem adotada e da organização dos
elementos da interface.

A análise dos aplicativos similares revelou uma linguagem visual comum entre as
soluções existentes, caracterizada por interfaces lúdicas, linguagem informal
e, em alguns casos, uso intenso de cores. Paralelamente, os resultados do
questionário indicaram que 37,5\% dos participantes preferem uma comunicação
acolhedora e empática, enquanto 25\% demonstraram preferência por uma linguagem
neutra e objetiva.

O design, no contexto de desenvolvimento de software é um elemento fundamental,
considerando que influencia diretamente a experiência do usuário, na percepção
e assimilação das informações.

Diferentemente da maioria dos aplicativos analisados, optou-se por uma
identidade visual menos lúdica e mais sóbria, reduzindo o uso de ilustrações e
cores saturadas em favor de uma linguagem mais discreta e orientada ao
conteúdo.

\section{Logotipo e identidade visual}

O logotipo do sistema foi concebido para representar o próprio nome da
aplicação. O termo Nexo remete à ideia de conexão, relação e convergência entre
diferentes elementos. Esse conceito dialoga diretamente com a proposta do
projeto, fundamentada na metodologia de Avaliação Ecológica Momentânea (EMA).
Embora seja cada vez mais comum monitorar aspectos isolados da rotina, como
qualidade do sono, nível de atividade física ou uso de medicamentos, essas
informações frequentemente permanecem fragmentadas, dificultando a compreensão
das relações existentes entre elas.

Nesse contexto, o nome Nexo representa a proposta central do sistema: integrar
essas informações e auxiliar o usuário na construção de uma narrativa coerente
sobre seu estado emocional. Em vez de tratar humor, emoções, estresse, energia,
ansiedade, atividades, psicofármacos e consultas como registros independentes,
o aplicativo busca evidenciar como esses fatores se relacionam ao longo do
tempo, favorecendo a identificação de padrões e possíveis associações.

Assim, tanto o nome quanto o logotipo procuram transmitir a ideia de
convergência, de ``dar a liga'', representando o papel do sistema em conectar
informações dispersas e transformá-las em conhecimento útil para o processo de
autorreflexão.

\begin{figure}[h]
\centering
\caption{Logotipo do sistema}
\label{fig:logo}

\includegraphics[width=0.9\linewidth]{logo.png}
\source{O autor}
\end{figure}

Para convergir com a pesquisa de tonalidade, e abraçar uma estética
``acolhedora'', o logo tenta fugir de parecer algo mecânico ao usar uma técnica
que simule pinceladas ao invés de linhas retas, simulando também diferenças ou
falhas na pressão utilizada por ter pequenas variações e espaços em branco no
logotipo.

\subsection{Hierarquia visual}

No contexto de interfaces digitais, a hierarquia visual corresponde à
estruturação de elementos de uma tela ou página de forma a narrar e orientar o
usuário com os elementos que são mais significativos e que demandam mais
atenção.

Por meio de repetições de padrões de agrupamento de informações, contraste,
tamanhos, espaçamentos, busca-se reduzir o esforço necessário para encontrar
elementos importantes e quais ações necessitam de maior visibilidade. Segundo
\citet{Lidwell_2011} uma hierarquia visual bem definida, ajuda o usuário a
interpretar rápidamente a interface, reduzindo a carga cognitiva e tornando a
navegação mais intuitiva.

Portanto, com esse foco em mente, produziu-se uma lista pequena de componentes
com visual semelhante, reduzindo o contraste entre esses mesmo elementos, o que
ajuda justamente a gerar uma sensação de consistência e previsibilidade entre
telas e elementos gráficos.

A estrutura da aplicação pode ser resumida em cinco componentes fundamentais:
cartões (\textit{cards}), botões de ação, \textit{chips} para categorização de
informações, campos de formulário e indicadores em escala (\textit{sliders}).
Esses componentes aparecem de maneira recorrente em praticamente todas as telas
do sistema, permitindo que o usuário reconheça rapidamente sua função sem
necessidade de reaprendizado.

Portanto, a organização de telas geralmente se dá nesses mesmos parâmetros, em
que toda tela é dividida em seções, compostas de um cabeçalho, seguido dos
respectivos elementos de interação. Posteriormente, a utilização de cartões,
ajuda a reforçar esse foco visual, agrupando elementos que são interativos e/ou
que necessitam interação e/ou atenção.

\subsection{Paleta de Cores}

Como discutido anteriormente, a definição da paleta de cores buscou equilibrar
dois objetivos identificados durante a fase de pesquisa, transmitir uma
identidade visual acolhedora, porém evitando excessos de cores e informações
lúdicas. Portanto, para buscar esses pilares, optou-se por usar uma escolha de
cores mais neutras e sóbrias, focando em tons de branco, cinza e preto, de
forma a permitir que o conteúdo regstrado pelo usuário seja o maior foco de
destaque.

\begin{figure}[h]
  \centering
  \caption{Paleta de cores}
  \label{fig:palette}
  \includegraphics[width=10cm, height=8cm, keepaspectratio]{pallete-small.png}
  \source{O autor}
\end{figure}

% Hack
\newpage

Contudo, para elementos que necessitam de uma maior atenção, ou elementos de
interação, como botões de ações primárias, adotou-se uma escolha mais rádical,
um tom de verde forte, que busca criar exatamente a sensação de destaque
imediato. Além disso, o verde geralmente está associado como uma cor que
transmite uma sensação de vitalidade, força, natureza, bem-estar e afins. Como
mencionado anteriormente, limitou-se a utilização para indicadores, componentes
de interação e elementos que representam progresso ou confirmação.

Por fim, foi escolhido outras cores pasteis para ajudar a destacar elementos
como sentimento e emoções, indicadores para cada tipo de registro gerado no
sistema, de forma mínima e não intrusiva.

\subsection{Tipografia}

A tipografia constitui um dos principais elementos responsáveis pela
legibilidade em interfaces digitais. Segundo o Material Design
\cite{Google_2021}, a escolha adequada da fonte influencia diretamente a
velocidade de leitura, a organização da informação e a percepção geral de
qualidade da interface. Em aplicações móveis, onde o espaço disponível é
reduzido, a clareza tipográfica torna-se ainda mais importante.

Considerando esses aspectos, optou-se pela utilização da fonte Inter,
desenvolvida por Rasmus Andersson especificamente para ambientes digitais. A
família tipográfica foi projetada para manter elevada legibilidade em
diferentes tamanhos de tela, apresentando boa diferenciação entre caracteres e
excelente renderização em sistemas operacionais modernos \cite{Andersson_2016}.

A hierarquia tipográfica do sistema foi construída a partir de poucas variações
de tamanho e peso da fonte. Títulos utilizam maior peso tipográfico para
estabelecer prioridade visual, enquanto textos informativos e descrições
empregam pesos mais leves, mantendo consistência em toda a aplicação. Essa
padronização reduz a quantidade de estímulos visuais simultâneos e facilita a
identificação das informações mais relevantes em cada tela.

Além da legibilidade, a escolha da Inter também contribui para reforçar a
identidade moderna e funcional proposta pelo aplicativo. Por ser amplamente
utilizada em plataformas digitais e ferramentas de interface, como o Figma, a
fonte transmite familiaridade ao usuário e mantém consistência com os padrões
contemporâneos de design de interfaces.
